% !TEX program = lualatex
% Auto-generated from syllabus — 06: PCの設定1：基本設定とネットワーク設定

\documentclass[aspectratio=43,professionalfonts,handout,10pt]{beamer}
\usepackage{beamer_template}

\newcommand{\LectureNum}{06}
\newcommand{\LectureTitle}{PCの設定1：基本設定とネットワーク設定}
\newcommand{\CourseName}{情報リテラシー}
\newcommand{\TextPages}{-}

\begin{document}
% --- Slide 1: title ---

\begin{frame}{\CourseName\quad 第\LectureNum 回「\LectureTitle 」}
  {\small \TermName\quad \TeacherName\quad ／\quad 教科書 \TextPages}

  \vfill

  \begin{block}{今日の流れ}
    \begin{enumerate}
      \item PCの初期セットアップ（アカウント・PIN・デバイス名）
      \item Windowsの基本設定
      \item ネットワーク接続（有線LAN→Wi-Fi）
    \end{enumerate}
  \end{block}
  \vfill

  \begin{block}{今日の目標}
    \begin{itemize}
      \item 各自のラップトップPCを設定し，本学のネットワークに接続することができる。
    \end{itemize}
  \end{block}
  \vfill

  \begin{exampleblock}{キーワード}
    \small \textbf{OS} ／ \textbf{Microsoftアカウント} ／ \textbf{Wi-Fi} ／ \textbf{Windows Update}
  \end{exampleblock}
\end{frame}

% --- Slide 2: 作業の流れ ---

\begin{frame}{本日の作業の流れ}
  \begin{block}{ステップ}
    \begin{itemize}
      \item ① PC起動・初期セットアップ（アカウント・PIN・デバイス名）
      \item ② Windowsの基本設定（時刻・更新・電源・ディスプレイ）
      \item ③ ネットワーク接続（有線LAN確認→Wi-Fi接続）
    \end{itemize}
  \end{block}

  \vfill

  \begin{noticebox}[注意]
    わからないときは自己判断せず，担当教員・技術職員に相談すること。
  \end{noticebox}
\end{frame}

% --- Slide 3: STEP 1 ---

\begin{frame}{STEP 1：Windowsの初期セットアップ}
  {\small PCを初めて起動し，Windowsのセットアップを完了する。}

  \vfill

  \begin{block}{手順}
    \begin{enumerate}
      \item ACアダプタとUSB有線LANアダプタを接続し，電源を入れる
      \item 画面の指示に従って言語を「日本語」，地域を「日本」，キーボードを「Microsoft IME」に設定する
      \item 第2回で作成した個人Microsoftアカウントでサインインする
      \item PINの設定を求められたら，4桁以上の数字を設定する
      \item デバイス名の入力画面で，英語の名前を設定する
      \item デスクトップが表示されたら次のSTEPに進む
    \end{enumerate}
  \end{block}
\end{frame}

% --- Slide 4: STEP 2 ---

\begin{frame}{STEP 2：Windowsの基本設定}
  {\small 快適・安全に使うための基本設定を行う。}

  \vfill

  \begin{block}{手順}
    \begin{enumerate}
      \item 「設定」→「時刻と言語」→日付・時刻が正しいか確認する
      \item 「設定」→「Windows Update」→更新プログラムを確認・インストールする
      \item 「設定」→「電源とバッテリー」→スリープまでの時間を15分に設定する
      \item 「設定」→「ディスプレイ」→解像度・スケールを確認する
    \end{enumerate}
  \end{block}
\end{frame}

% --- Slide 5: STEP 3 ---

\begin{frame}{STEP 3：学内ネットワークへの接続}
  {\small 有線LANでインターネット接続を確認した後，Wi-Fiにも接続する。}

  \vfill

  \begin{block}{手順}
    \begin{enumerate}
      \item ブラウザからインターネットにアクセスできることを確認する（有線LAN）
      \item LANケーブルを外し，タスクバーのWi-Fiアイコンから学内SSIDを選択する
      \item 配布されたユーザID・パスワードを入力して接続する
      \item Wi-Fiでもブラウザからインターネットにアクセスできることを確認する
    \end{enumerate}
  \end{block}
\end{frame}

% --- Slide 6: トラブル対応 ---

\begin{frame}{トラブル発生時の対処}
  \begin{noticebox}[よくあるトラブルと対処法]
    \begin{itemize}
      \item Wi-Fiに接続できない → SSIDとパスワードを再確認する
      \item Webページが開けない → ブラウザを変えて試す（Edge・Chrome等）
      \item セットアップが途中で止まる → 電源を切らずに担当教員・技術職員に報告する
      \item パスワードを忘れた → 担当教員・技術職員に申し出る（自己解決しようとしない）
    \end{itemize}
  \end{noticebox}

  \vfill

  \begin{exampleblock}{ポイント}
    機器の故障と操作ミスは見た目が似ている。まず担当教員・技術職員に相談すること。
  \end{exampleblock}
\end{frame}

% --- チェックリスト ---

\begin{frame}{まとめ・チェックリスト}
  \begin{block}{自己チェック（できたら $\checkmark$ をつけよう）}
    \begin{itemize}
      \item[$\square$] Windowsの初期セットアップを完了した（デバイス名を英語で設定）
      \item[$\square$] 基本設定（時刻・更新・電源・ディスプレイ）を確認した
      \item[$\square$] 有線LANでインターネットに接続できた
      \item[$\square$] 学内Wi-Fiに接続できた
    \end{itemize}
  \end{block}

\end{frame}

\end{document}
